\section{Límites y perspectivas de la construcción de aplicaciones de IA}

\subsection{Límites actuales}

\begin{warningbox}{``Everything is awesome until it breaks''}
El mayor reto de la construcción de aplicaciones de IA: cuando todo funciona bien la experiencia es excelente, pero en cuanto algo falla se vuelve al punto de partida: el desarrollador todavía necesita entender el código subyacente para depurarlo y corregirlo.
\end{warningbox}

\begin{itemize}
    \item \textbf{El prompt es solo una sugerencia}: no todo usuario lo entiende así
    \item \textbf{Seguridad}: históricamente han surgido problemas de seguridad
    \item \textbf{Homogeneización}: ¿cómo evitar que todas las aplicaciones generadas por IA se vean idénticas?
    \item \textbf{Límite de complejidad}: ¿qué tan complejas pueden llegar a ser estas aplicaciones?
\end{itemize}

\subsection{Impacto real}

Las herramientas de construcción de aplicaciones de IA están teniendo impacto en tres niveles:

\begin{itemize}
    \item \textbf{Startup Velocity}: las empresas emergentes lanzan un MVP en días en lugar de meses
    \item \textbf{Enterprise Modernization}: las organizaciones grandes validan primero con prototipos de IA antes de invertir recursos de ingeniería
    \item \textbf{Agency Efficiency}: las agencias de diseño entregan prototipos interactivos horas después de la reunión de arranque
\end{itemize}

\subsection{Criterios del equipo para elegir un App Builder}

Si el contenido del curso se traduce en una decisión de ingeniería, se pueden usar tres preguntas para evaluar si un App Builder de IA en verdad es adecuado para el equipo:

\begin{enumerate}
    \item \textbf{¿Permite iterar de manera continua y no solo generar una vez?} Una herramienta que en verdad aporte valor debe admitir múltiples rondas de modificación, verificación y reversión, en lugar de terminar tras una sola generación.
    \item \textbf{¿Respeta la pila tecnológica existente?} El equipo necesita una herramienta que produzca código que pueda integrarse al React/Next.js/sistema de diseño ya existente, no un prototipo aislado.
    \item \textbf{¿Permite bajar al nivel del código subyacente con rapidez cuando algo falla?} Si la herramienta se vuelve completamente indepurable en cuanto falla, solo sirve como generador de demostraciones, no como herramienta de producción.
\end{enumerate}

\begin{warningbox}{La velocidad del prototipo no es el único indicador}
Los App Builder de IA suelen venderse con el atractivo de «generar una aplicación en minutos», pero lo que en verdad determina el valor a largo plazo es el costo de modificación, el costo de depuración y la compatibilidad con los sistemas de ingeniería ya existentes. Una generación rápida con un costo de mantenimiento posterior alto termina generando una nueva deuda técnica.
\end{warningbox}

\subsection{Resumen de la sección}

Las herramientas de construcción de aplicaciones de IA han reducido de manera notable el tiempo entre la idea y el prototipo, pero el problema de que «en cuanto algo falla se vuelve al punto de partida» nos recuerda que entender la tecnología subyacente sigue siendo indispensable.

